\centerline{\bf \Huge Показательный листочек}
\centerline{\bf \Huge }

\begin{flushright}
        \scriptsize{}
\end{flushright}


\problem{1} 
Докажите, что если $p$ --- простое, то у числа $2^p-1$ все простые делители больше $p$.

\problem{2} 
Докажите, что если $p$ --- простое, то любой простой делитель числа $2^p-1$ имеет вид $2kp+1$ для некоторого натурального $k$.

\problem{3}
Докажите, что все простые делители чисел Ферма $2^{2^n}+1$ имеют вид $2^{n+1}\cdot x+1$. (Именно так Эйлер нашел простой делитель у числа $2^{2^5}+1$.)

\problem{4} 
Докажите, что для любого натурального $n$ число $2^n-1$ не делится на $n$.

\problem{5} 
Существует ли набор натуральных чисел $(n_1,\ldots,n_k)$, такой, что $n_i\mid 2^{n_{i+1}}-1$ для всех $i=1$, \ldots, $k$ (мы считаем, что $n_{k+1}=n_1$)?

\problem{6} 
Найдите все тройки простых чисел $(p,q,r)$, такие, что $p \mid q^r+1$, $q \mid r^p+1$, $r \mid p^q+1$.

\problem{7} 
Найти все пары простых чисел $(p;q)$, таких, что $pq \mid (5^p-2^p)(5^q-2^q)$.

\problem{8} 
Найти все пары простых чисел $p$, $q$, таких, что $pq \mid 2^p+2^q$.

\problem{9} 
Найдите все натуральные $n$, для которых числа $n$ и $2^n + 1$ имеют один и тот же набор простых делителей.

%\textbf{( b )} Найдите все натуральные $n$, для которых суммы
%простых делителей чисел $n$ и $2^n + 1$ равны.

%\task Найдите все натуральные $n$, такие, что число $\cfrac{2^n +
%1}{n^2}$ является целым.
